\chapter{\textit{Main Contribution}}
\label{ch:mainmatter}

This chapter contains the main theoretical (and possibly technical) contribution of the thesis.
It should present a coherent body of new results, organized into sections.

For a theory-focused thesis: introduce new notions, state lemmas/theorems,
provide proofs, and illustrate results with examples.
For an implementation-focused thesis: additionally describe the algorithm,
key design decisions, and data structures.

The opening paragraph of this chapter should cover the following points:
\begin{enumerate}
    \item State what new objects, criteria, or techniques this chapter introduces.
    \item Introduce a running example that will be used throughout to illustrate
          each new definition and result.
    \item Outline the chapter structure by referencing the main sections.
\end{enumerate}

%-----------------------------------------------------------------------
\section{[\textit{Central Definitions / New Framework}]}\label{sec:main-definitions}
%-----------------------------------------------------------------------

% Introduce the key new notions.
% Follow the style of \cref{ch:preliminaries}: state each definition formally,
% add a brief intuitive explanation, and accompany it with at least one example.
% Move from simple to complex.

For self-introduced notions it is even more important to explain everything
and use multiple examples to illustrate everything.
Use one running example (\Cref{ex01}) that can be used through the whole section 
and maybe some additional small examples to highlight special cases.

Each result in this section should be presented as follows:
\begin{enumerate}
    \item State and prove all auxiliary lemmas before the main theorem.
    \item State the main theorem.  Typical statements are:
          \begin{itemize}
              \item \textbf{Soundness:} If the technique returns \textsc{Yes},
                    then the property holds.
              \item \textbf{Completeness:} If the property holds, the technique
                    returns \textsc{Yes}.
              \item \textbf{Equivalence:} The new criterion holds if and only if
                    an existing criterion holds.
          \end{itemize}
    \item Provide a complete proof, structured into numbered cases for long
          arguments.  Reference definitions and earlier results by label.
    \item Apply the theorem to the running example from \Cref{ex01},
          walking through every step explicitly.
\end{enumerate}

\section{Example Section}

Lorem ipsum dolor sit amet, consectetur adipiscing elit \cite{Doe2024Example}. \todo{Rephrase this.}
\fdbk{Further details required.}

% In the following, you can see a concrete example as if we would define our own new type of DTMCs

A My discrete-time Markov chain (MDTMC) models a system that moves between states in discrete steps, where the probability of the next state depends only on the current state and one state in the state space is absorbing.

\begin{definition}[My Discrete-time Markov chain (MDTMC)]
A \emph{My discrete-time Markov chain (MDTMC)} is a stochastic process $(X_n)_{n\in\mathbb{N}}$ on a (countable) state space $S$ such that for all $n$ and $s_0,\dots,s_{n},s'\in S$:
$$
\Pr(X_{n+1}=s' \mid X_n=s_n, X_{n-1}=s_{n-1},\dots,X_0=s_0)
=\Pr(X_{n+1}=s' \mid X_n=s_n).
$$

Moreover, there exists a transition matrix $P=(p_{s,s'})_{s,s'\in S}$ with
$p_{s,s'}=\Pr(X_{n+1}=s'\mid X_n=s)$ (independent of $n$) and
$\sum_{s'\in S} p_{s,s'}=1$ for all $s\in S$.
We assume that the underlying DTMC has exactly one absorbing state, i.e.\ there exists at most one $s\in S$ with
$p_{s,s}=1$.
\end{definition}

\begin{example}[A small MDTMC Example]
\label{ex01}
Let $S=\{0,1,2\}$.
From state $0$, the chain stays in $0$ with probability $0.7$ and moves to $1$ with probability $0.3$.
From state $1$, it moves to $2$ with probability $0.5$ and to $0$ with probability $0.5$.
State $2$ is absorbing.
\Cref{tab01} shows the corresponding transition matrix.
\end{example}

\begin{table}[H]
\centering
\caption{Transition matrix $P$ of the example MDTMC from \Cref{ex01}.}
\begin{tabular}{c|ccc}
$P$ & 0 & 1 & 2\\ \hline
0 & 0.7 & 0.3 & 0.0\\
1 & 0.5 & 0.0 & 0.5\\
2 & 0.0 & 0.0 & 1.0\\
\end{tabular}
\label{tab01}
\end{table}

\begin{theorem}[Key Theorem]
    \label{keythm}
    Every MDTMC has an absorving state.
\end{theorem}

\begin{proof}
    The proof is left as an exercise for the reader.
\end{proof}